%% ============================================================
%%  二次元気液混合流体の数値計算法 ── 教科書的総合解説
%%  Conservative Level Set法 + 結合コンパクト差分法(CCD) + 界面適合格子
%%  + Projection法 による完全定式化
%%
%%  XeLaTeX でコンパイル
%% ============================================================
\documentclass[12pt, a4paper, titlepage, xelatex, ja=standard]{bxjsarticle}

%% ── フォント設定（XeLaTeX + xeCJK）──────────────────────────────
\usepackage{fontspec}
\usepackage{xeCJK}

\setmainfont{Times New Roman}
\setCJKmainfont{Hiragino Mincho ProN}

%% ── 数学・図表パッケージ ──────────────────────────────────────
\usepackage{amsmath, amssymb, amsthm, mathtools}
\usepackage{bm}
\usepackage{booktabs, array, multirow, tabularx}
\usepackage{graphicx}
\usepackage{xcolor}
\usepackage{tcolorbox}
\tcbuselibrary{skins, breakable, theorems}
\usepackage{tikz}
\usetikzlibrary{arrows.meta, positioning, shapes.geometric, backgrounds, calc, decorations.pathreplacing, matrix, patterns}
\usepackage{geometry}
\usepackage{fancyhdr}
\usepackage{titlesec}
\usepackage{hyperref}
\usepackage{cleveref}
\usepackage{enumitem}
\usepackage{caption}
\usepackage{subcaption}
\usepackage{microtype}
\usepackage{nicefrac}
\usepackage{ragged2e}

%% ── ページレイアウト ──────────────────────────────────────────
\geometry{top=25mm, bottom=28mm, left=26mm, right=26mm, headheight=16pt}
\emergencystretch=2em
\setlength{\hfuzz}{3pt}

%% ── カラー定義 ────────────────────────────────────────────────
\definecolor{navy}{RGB}{0,40,90}
\definecolor{blue2}{RGB}{0,80,160}
\definecolor{midblue}{RGB}{0,90,170}
\definecolor{lightblue}{RGB}{230,242,255}
\definecolor{tblue}{RGB}{218,234,255}
\definecolor{tgreen}{RGB}{215,255,225}
\definecolor{tyellow}{RGB}{255,250,210}
\definecolor{tred}{RGB}{255,225,220}
\definecolor{darkgray}{RGB}{60,60,60}
\definecolor{boxgray}{RGB}{245,245,245}
\definecolor{green4}{RGB}{0,120,60}
\definecolor{red2}{RGB}{160,0,0}
\definecolor{orange2}{RGB}{200,90,0}
\definecolor{torange}{RGB}{255,238,210}
\definecolor{teal2}{RGB}{0,110,110}
\definecolor{tteal}{RGB}{210,245,245}
\definecolor{purple2}{RGB}{100,0,140}
\definecolor{tpurple}{RGB}{240,220,255}
\definecolor{brown2}{RGB}{120,60,0}
\definecolor{tbrown}{RGB}{250,235,215}

%% ── tcolorbox スタイル ────────────────────────────────────────
\tcbset{
  mybox/.style={
    enhanced, breakable,
    colback=tblue, colframe=blue2,
    toprule=0.8pt, bottomrule=0.8pt, leftrule=0.8pt, rightrule=0.8pt,
    fonttitle=\bfseries\small\color{white},
    colbacktitle=navy,
    attach boxed title to top left={yshift=-2mm, xshift=6mm},
    boxed title style={sharp corners},
    sharp corners=south,
  },
  derivbox/.style={
    enhanced, breakable,
    colback=boxgray, colframe=navy!40,
    toprule=0.5pt, bottomrule=0.5pt, leftrule=3pt, rightrule=0.5pt,
    sharp corners,
  },
  resultbox/.style={
    enhanced, colback=tgreen, colframe=navy,
    toprule=1pt, bottomrule=1pt, leftrule=1pt, rightrule=1pt,
    sharp corners,
  },
  warnbox/.style={
    enhanced, colback=tyellow, colframe=red2!70,
    toprule=0.6pt, bottomrule=0.6pt, leftrule=3pt, rightrule=0.6pt,
    sharp corners,
  },
  algbox/.style={
    enhanced, breakable, colback=white, colframe=navy,
    toprule=1pt, bottomrule=1pt, leftrule=1pt, rightrule=1pt,
    sharp corners,
    fonttitle=\bfseries\small\color{white}, colbacktitle=navy,
    attach boxed title to top left={yshift=-2mm, xshift=6mm},
    boxed title style={sharp corners},
  },
  defbox/.style={
    enhanced, colback=boxgray, colframe=darkgray!50,
    toprule=0.6pt, bottomrule=0.6pt, leftrule=0.6pt, rightrule=0.6pt,
    sharp corners,
  },
}

%% ── ヘッダー/フッター ─────────────────────────────────────────
\pagestyle{fancy}
\fancyhf{}
\fancyhead[L]{\small\color{navy}二次元気液混合流体の数値計算法}
\fancyhead[R]{\small\color{navy!70}Conservative Level Set~+~CCD}
\fancyfoot[C]{\small\thepage}
\renewcommand{\headrulewidth}{0.4pt}
\renewcommand{\footrulewidth}{0pt}

%% ── セクション書式 ────────────────────────────────────────────
\titleformat{\section}{\large\bfseries\color{navy}}{\thesection.}{0.8em}{}[\color{blue2}\rule{\linewidth}{0.8pt}]
\titleformat{\subsection}{\normalsize\bfseries\color{navy!80}}{\thesubsection}{0.8em}{}
\titleformat{\subsubsection}{\normalsize\bfseries\color{navy!60}}{\thesubsubsection}{0.8em}{}

%% ── 数学マクロ ────────────────────────────────────────────────
\newcommand{\pd}[2]{\dfrac{\partial #1}{\partial #2}}
\newcommand{\pdd}[2]{\dfrac{\partial^2 #1}{\partial #2^2}}
\newcommand{\pdm}[3]{\dfrac{\partial^2 #1}{\partial #2\,\partial #3}}
\newcommand{\bu}{\bm{u}}
\newcommand{\bg}{\bm{g}}
\newcommand{\bnabla}{\bm{\nabla}}
\newcommand{\He}{H_\varepsilon}
\newcommand{\de}{\delta_\varepsilon}
\newcommand{\Ord}[1]{\mathcal{O}(#1)}
\newcommand{\half}{\tfrac{1}{2}}
\newcommand{\fp}[1]{f^{(1)}_{#1}}
\newcommand{\fpp}[1]{f^{(2)}_{#1}}

%% theorem 環境
\newtheorem{remark}{Remark}
\newtheorem{definition}{定義}
\newtheorem{theorem}{定理}

%% ── ハイパーリンク ────────────────────────────────────────────
\hypersetup{
  colorlinks=true, linkcolor=navy, citecolor=navy, urlcolor=blue2,
  pdftitle={Two-Dimensional Gas-Liquid Two-Phase Flow Numerical Methods},
}

%% ============================================================
\begin{document}
\RaggedRight
%% ============================================================

%% ── 表紙 ──────────────────────────────────────────────────────
\begin{titlepage}
\begin{center}
\vspace*{12mm}
{\color{navy}\rule{\linewidth}{2pt}}\\[8pt]
{\LARGE\bfseries 二次元気液混合流体の数値計算法}\\[6pt]
{\large\bfseries ── 基礎理論から完全アルゴリズムまでの教科書的解説 ──}\\[4pt]
{\color{navy}\rule{\linewidth}{2pt}}\\[10pt]

{\large
\textbf{Numerical Methods for Two-Dimensional Gas--Liquid Two-Phase Flows}\\[4pt]
\textit{From Basic Theory to Complete Algorithm}\\[3pt]
\textit{Coupled Compact Difference $\cdot$ Conservative Level Set $\cdot$ Projection Method}
}

\vspace{18mm}
\begin{tcolorbox}[defbox, width=0.85\textwidth]
\small\centering
\begin{tabular}{rl}
界面追跡法： & 質量保存型 Conservative Level Set 法 (Olsson \& Kreiss)\\[2pt]
空間離散化： & 結合コンパクト差分法（CCD 法）$\Ord{h^6}$\\
                  & Chu \& Fan, \textit{J. Comput. Phys.}, 1998\\[2pt]
移流安定化： & 5次精度 WENO スキーム\\[2pt]
格子構造：   & コロケート格子 + Rhie--Chow 補正\\[2pt]
界面適合格子： & $\psi$ 連動グリッド密度関数による自動高解像度化\\[2pt]
時間積分：   & 3次 TVD Runge--Kutta （対流・Level Set）\\
                  & Crank--Nicolson （粘性項，半陰的）\\[2pt]
圧力求解：   & 変密度 Chorin Projection + BiCGSTAB\\[2pt]
表面張力：   & CSF 法（連続体表面力法）
\end{tabular}
\end{tcolorbox}

\vspace{10mm}
{\small\color{darkgray}
本稿は、気液二相流における質量保存の本質的課題を解決する\\
Conservative Level Set 法と 5次精度 WENO スキームの導入を中心に、\\
高精度 CCD 法を組み合わせた完全な定式化を解説したものである。
}
\end{center}

\vfill
{\color{navy}\rule{\linewidth}{0.8pt}}\\[2pt]
{\small\color{darkgray}対象読者：大学院生・研究者（計算流体力学の基礎知識を前提とする）}
\end{titlepage}

%% ── 目次 ──────────────────────────────────────────────────────
\tableofcontents
\newpage

%% ============================================================
\section{はじめに}
\label{sec:intro}
%% ============================================================

\subsection{気液二相流の数値計算における課題}

気液二相流とは，気相（ガス）と液相（液体）が共存し，
その間に明確な界面を形成しながら運動する流体現象である．
気泡塔，核沸騰，波浪，噴霧燃焼など，工学・自然界に広く現れる流体現象である．

この現象を数値的に扱う際には，以下の4つの本質的困難が存在する：

\begin{enumerate}
  \item \textbf{界面の捕捉}：密度・粘性が界面を挟んで不連続に変化する
        （$\rho_l/\rho_g \approx 1000$，$\mu_l/\mu_g \approx 100$）ため，
        界面位置を高精度かつ安定に追跡しなければならない．
  \item \textbf{表面張力の正確な評価}：表面張力は界面の曲率（$\kappa$）に比例する体積力である．
        曲率の計算には二階偏微分が必要となり，数値誤差が増幅されやすい．
  \item \textbf{非圧縮条件の高精度維持}：速度場の発散をゼロに保つ
        圧力場を毎ステップ正確に求解する必要がある．
  \item \textbf{密度成層の安定処理}：重力と密度差による静水圧分布を
        正確に扱わないと寄生流れ（parasitic currents）が発生する．
\end{enumerate}

\subsection{本稿のアプローチ}

上記の困難を克服するため，本稿では以下の手法を組み合わせる
（図\ref{fig:ns_solvers}および図\ref{fig:timestepping}参照）：

\begin{itemize}
  \item \textbf{界面追跡：} Conservative Level Set 法 \cite{Olsson2005} 
        を採用する。従来の Level Set 法と異なり、移流方程式を保存形
        $\partial \psi/\partial t + \bnabla \cdot (\psi \bu) = 0$ で解くことで、
        長時間の計算においても全質量（体積）が厳密に保存される。
  \item \textbf{高精度空間微分：} 結合コンパクト差分法（CCD 法，Chu \& Fan \cite{ChuFan1998}）を用い，
        一次微分・二次微分・混合偏微分をすべて $\Ord{h^6}$ 精度で統一的に評価する．
        本稿では係数をテイラー展開から厳密に導出する．
  \item \textbf{界面適合格子：} Conservative Level Set 関数 $\psi$ に連動したグリッド密度関数 $\omega(\psi)$ により，
        界面近傍のみを自動的に高解像度化する非一様格子を生成する．
  \item \textbf{格子構造：} チェッカーボード不安定を防ぐため
        コロケート格子を採用し，Rhie--Chow 型流束補正を適用する．
  \item \textbf{時間積分：} 対流項と Level Set 移流には
        3次 TVD Runge--Kutta を，粘性項には Crank--Nicolson 半陰的処理を用いる．
  \item \textbf{圧力求解：} 変密度 Chorin Projection 法に基づいて圧力 Poisson 方程式を導き，
        CCD で差分化した後 BiCGSTAB で反復求解する．
\end{itemize}

%%% ── 図1：NS各項とソルバーの対応 ────────────────────────────
\begin{figure}[h]
\centering
\begin{tikzpicture}[
  node distance=4mm and 5mm,
  nsterm/.style={draw=navy!70, fill=lightblue, rounded corners=4pt,
                 font=\small\bfseries, inner sep=6pt,
                 minimum width=26mm, minimum height=10mm, align=center},
  solver/.style={draw=#1!70, fill=#1!10, rounded corners=4pt,
                 font=\scriptsize, inner sep=5pt,
                 minimum width=26mm, align=center},
  solver/.default=navy,
  subarr/.style={-Stealth, semithick, darkgray!60},
]

\node[font=\small\bfseries, text=navy] (hd1) at (0,0)
  {NS 方程式の各項とソルバーの対応};

%% ── NS 5項ボックス ──────────────────────────────────────────
\node[nsterm, below=7mm of hd1, xshift=-57mm] (T1) {対流項};
\node[nsterm, right=4mm of T1]  (T2) {圧力項};
\node[nsterm, right=4mm of T2]  (T3) {粘性項};
\node[nsterm, right=4mm of T3]  (T4) {重力項};
\node[nsterm, right=4mm of T4]  (T5) {表面張力項};

%% ── 各ソルバーボックス ──────────────────────────────────────
\node[solver=orange2, below=5mm of T1] (S1)
  {CCD（$\Ord{h^6}$）\\TVD-RK3 + WENO5\\\陽的};
\node[solver=midblue, below=5mm of T2] (S2)
  {CCD $\to\bnabla p$\\FVM-PPE\\Projection法};
\node[solver=teal2, below=5mm of T3] (S3)
  {CCD（$\Ord{h^6}$）\\Crank-Nicolson\\半陰的};
\node[solver=darkgray, below=5mm of T4] (S4)
  {陽的\\(定数項)};
\node[solver=red2, below=5mm of T5] (S5)
  {Cons. Level Set $\to\kappa$\\CCD\\半陰的};

\foreach \t/\s in {T1/S1, T2/S2, T3/S3, T4/S4, T5/S5}
  \draw[subarr] (\t) -- (\s);

\end{tikzpicture}
\caption{NS 方程式の各項と対応ソルバーの概要．
対流項は CCD + 5次精度 WENO スキームを組み合わせた TVD-RK3 で陽的に処理し，
圧力項は FVM ベース PPE を Projection 法で求解する．
粘性項は CCD + Crank--Nicolson で半陰的に扱い，
表面張力項は Conservative Level Set 関数から曲率を CCD で計算して評価する．}
\label{fig:ns_solvers}
\end{figure}

%%% ── 図2：タイムステッピングフロー ──────────────────────────
\begin{figure}[h]
\centering
\begin{tikzpicture}[
  node distance=4mm and 5mm,
  phasebox/.style={draw=#1!80, fill=#1!8, rounded corners=4pt,
                   font=\small, inner sep=6pt,
                   text width=115mm, align=left},
  phasebox/.default=navy,
  statebox/.style={draw=navy!60, fill=lightblue!50, rounded corners=4pt,
                   font=\small, inner sep=6pt,
                   text width=115mm, align=center},
  arr/.style={-Stealth, thick, navy},
]

\node[font=\small\bfseries, text=navy] (hd2) at (0,0)
  {タイムステッピングフロー（$t^n \to t^{n+1}$）};

\node[statebox, below=5mm of hd2] (In)
  {入力：\,$\psi^n$，$\bu^n$，$p^n$，$\tilde\rho^n$};

\node[phasebox=red2, below=3mm of In] (P1)
  {\textcolor{red2}{\textbf{①② Conservative Level Set 更新}}\enspace
   保存形移流（WENO5 + TVD-RK3）→ 再圧縮・正規化};

\node[phasebox=green4, below=2.5mm of P1] (P2)
  {\textcolor{green4}{\textbf{③④⑤ 格子・物性・曲率更新}}\enspace
   密度関数 $\omega(\psi^{n+1})$ → 物性値線形補間 → CCD で $\kappa^{n+1}$};

\node[phasebox=navy, below=2.5mm of P2] (P3)
  {\textcolor{navy}{\textbf{⑥ 運動量予測（Predictor）}}}\\\jp{対流（WENO5）・粘性（半陰的）・重力・表面張力 → 予測速度 $\bu^*$};

\node[phasebox=orange2, below=2.5mm of P3] (P4)
  {\textcolor{orange2}{\textbf{⑦ 圧力 Poisson 求解（FVM）}}}\enspace
   Rhie-Chow 補正 → 変密度 PPE → BiCGSTAB → $p^{n+1}$};

\node[phasebox=purple2, below=2.5mm of P4] (P5)
  {\textcolor{purple2}{\textbf{⑧⑨ 速度補正・収束確認}}}\enspace
   CCD $\bnabla p^{n+1}$ + Rhie-Chow 補正 → $\bu^{n+1}$；CFL更新};

\node[statebox, below=3mm of P5] (Out)
  {出力：\,$\psi^{n+1}$，$\bu^{n+1}$，$p^{n+1}$\quad$\to$\quad$t:=t+\Delta t$};

\foreach \a/\b in {In/P1, P1/P2, P2/P3, P3/P4, P4/P5, P5/Out}
  \draw[arr] (\a) -- (\b);

%% ループ矢印
\draw[arr, rounded corners=5pt]
  (Out.east) -- ++(7mm,0) -- ++(0,46mm)
  node[midway, right, font=\scriptsize, rotate=90, anchor=south,
       text=navy, yshift=0.5mm] {次ステップへ}
  -- (In.east);

\end{tikzpicture}
\caption{タイムステッピングフロー（$t^n \to t^{n+1}$）．
①〜⑨ は第\ref{sec:algorithm}章の完全アルゴリズムフローに対応する．
Conservative Level Set の更新によって質量の不変性を保証してから、
高精度な速度・圧力計算へと進む。}
\label{fig:timestepping}
\end{figure}

\subsection{本稿の構成}

第\ref{sec:governing}章で支配方程式（One-Fluid Navier--Stokes + Level Set）を導く．
第\ref{sec:levelset}章で Conservative Level Set 法の詳細（保存形移流・再圧縮・WENO）を述べる．
第\ref{sec:grid}章で界面適合非一様格子の構築法と座標変換を説明する．
第\ref{sec:CCD}章が本稿の数学的核心であり，CCD スキームの係数をテイラー展開から完全に導出する．
第\ref{sec:collocate}章でコロケート格子と Rhie--Chow 補正を述べる．
第\ref{sec:pressure}章で圧力 Poisson 方程式を導出し CCD で差分化する．
第\ref{sec:time}章で時間積分スキームを説明し，
第\ref{sec:algorithm}章で完全アルゴリズムをまとめる．


%% ============================================================
\section{支配方程式系の導出}
\label{sec:governing}
%% ============================================================

\subsection{One-Fluid（単一流体）定式化の概念}

気相・液相を別々の変数で扱う代わりに，
界面近傍で物性値を滑らかに補間した \textbf{「一つの流体」として扱う}定式化を One-Fluid 法という．
変密度・変粘性の単一 Navier--Stokes 方程式で全領域を記述できるため，
界面の陰関数表現（Level Set など）と自然に組み合わせられる．

\subsection{非圧縮条件（連続の式）}

気相・液相とも非圧縮性流体（$\rho = $const in each phase）を仮定する．
One-Fluid 定式化では，界面が移流によって運ばれるだけで質量が生成・消滅しないから，
速度場に対して以下の拘束条件が成立する：
\begin{equation}
  \bnabla \cdot \bu = 0
  \label{eq:continuity}
\end{equation}
この式は各タイムステップで圧力場を通じて強制的に満たされる（後述の Projection 法）．

\subsection{Navier--Stokes 方程式（One-Fluid + CSF）}

One-Fluid 定式化による運動方程式は，密度・粘性が空間的に変化することを許した
Navier--Stokes 方程式である：

\begin{tcolorbox}[mybox, title={One-Fluid Navier--Stokes 方程式（次元形）}]
\begin{equation}
  \rho(\psi)\!\left(\pd{\bu}{t} + \bu\cdot\bnabla\bu\right)
  = -\bnabla p
  + \bnabla\cdot\!\left[\mu(\psi)\!\left(\bnabla\bu + \bnabla\bu^T\right)\right]
  + \rho(\psi)\,\bg
  + \sigma\,\kappa\,\bnabla \psi
  \label{eq:NS_dim}
\end{equation}
\end{tcolorbox}

各項の物理的意味を表\ref{tab:NS_terms}に整理する．

\begin{table}[h]
\centering
\caption{Navier--Stokes 方程式の各項の意味}
\label{tab:NS_terms}
\begin{tabular}{@{}lll@{}}
\toprule
項 & 数式 & 物理的意味\\
\midrule
慣性力（左辺） & $\rho D\bu/Dt$ & 質量$\times$加速度（Newton 第2法則）\\
圧力勾配力 & $-\bnabla p$ & 等方圧縮力\\
粘性応力 & $\bnabla\cdot[\mu(\bnabla\bu+\bnabla\bu^T)]$ & 変密度対応の Stokes 応力テンソル\\
重力 & $\rho\,\bg$ & $\bg = -g\hat{\bm{z}}$（下向き重力）\\
表面張力（CSF法） & $\sigma\kappa\bnabla\psi$ & 界面の曲率を体積力に変換 \cite{Olsson2005}\\
\bottomrule
\end{tabular}
\end{table}

\subsubsection{CSF 法（連続体表面力法）の導出}

界面上の表面張力ベクトルを、Conservative Level Set 関数 $\psi$ の勾配を用いて体積力として記述する。
$\psi$ が $0 \to 1$ へ滑らかに変化するステップ関数の場合、その勾配 $\bnabla \psi$ は界面法線方向に
デルタ関数的な重みを持つ。
したがって、曲率 $\kappa$ を用いた表面張力項は：
\begin{equation}
  \bm{F}_{\mathrm{CSF}} = \sigma\kappa\bnabla\psi
\end{equation}
と簡潔に表記される。これは従来の $\sigma\kappa\delta(\phi)\bnabla\phi$ と同等の効果を持つが、
$\psi$ 自身が滑らかな階段関数であるため、追加のデルタ関数近似を必要としない利点がある。

\subsection{界面曲率の計算式}

平均曲率 $\kappa$ は Level Set 関数の発散（単位法線ベクトルの発散）として定義される：
\begin{equation}
  \kappa(\psi) = -\bnabla\cdot\!\left(\frac{\bnabla\psi}{|\bnabla\psi|}\right)
  \label{eq:curvature_def}
\end{equation}
2次元の場合にこれを展開すると，$\psi_x = \partial\psi/\partial x$ 等と略記して：
\begin{tcolorbox}[mybox, title={2次元界面曲率（展開形）}]
\begin{equation}
  \kappa = -\frac{
    \psi_x^2\,\psi_{yy}
    - 2\psi_x\psi_y\,\psi_{xy}
    + \psi_y^2\,\psi_{xx}
  }{\bigl(\psi_x^2 + \psi_y^2\bigr)^{3/2}}
  \label{eq:curvature_2d}
\end{equation}
\end{tcolorbox}
この計算には $\psi_{xx}$，$\psi_{yy}$，$\psi_{xy}$ という二階偏微分が必要であり，
CCD 法を用いることで $\Ord{h^6}$ 精度の高精度評価が可能となる（第\ref{sec:CCD}章参照）．

\subsubsection{曲率計算の導出（式\eqref{eq:curvature_def}から式\eqref{eq:curvature_2d}）}

$\hat{\bm{n}} = (\psi_x/|\bnabla\psi|,\, \psi_y/|\bnabla\psi|)$ として発散を計算する：
\begin{align}
  \kappa
  &= -\left(\pd{}{x}\frac{\psi_x}{|\bnabla\psi|} + \pd{}{y}\frac{\psi_y}{|\bnabla\psi|}\right) \notag\\
  &= -\frac{\psi_{xx}|\bnabla\psi|^2 - \psi_x(\psi_x\psi_{xx}+\psi_y\psi_{xy})}{|\bnabla\psi|^3} -\frac{\psi_{yy}|\bnabla\psi|^2 - \psi_y(\psi_x\psi_{xy}+\psi_y\psi_{yy})}{|\bnabla\psi|^3} \notag\\
  &= -\frac{(\psi_x^2+\psi_y^2)(\psi_{xx}+\psi_{yy}) - (\psi_x^2\psi_{xx}+2\psi_x\psi_y\psi_{xy}+\psi_y^2\psi_{yy})}{|\bnabla\psi|^3} \notag\\
  &= -\frac{\psi_y^2\psi_{xx} - 2\psi_x\psi_y\psi_{xy} + \psi_x^2\psi_{yy}}{|\bnabla\psi|^3}
\end{align}

\subsection{無次元化}

参照量として液相密度 $\rho_l$，代表速度 $U$，代表長さ $L$ を選ぶ：
\begin{equation}
  \tilde{\bm{x}} = \frac{\bm{x}}{L},\quad \tilde{\bu} = \frac{\bu}{U},\quad \tilde{t} = \frac{Ut}{L},\quad \tilde{p} = \frac{p}{\rho_l U^2},\quad \tilde{\psi} = \psi \label{eq:nondim_vars}
\end{equation}
無次元パラメータ（チルダは省略）：
\begin{equation}
  Re = \frac{\rho_l U L}{\mu_l},\quad Fr = \frac{U}{\sqrt{gL}},\quad We = \frac{\rho_l U^2 L}{\sigma},\quad \hat{\rho} = \frac{\rho_g}{\rho_l},\quad \hat{\mu} = \frac{\mu_g}{\mu_l} \label{eq:nondim_params}
\end{equation}
無次元密度・粘性係数：
\begin{equation}
  \rho(\psi) = \hat{\rho} + (1-\hat{\rho})\psi,\quad \mu(\psi) = \hat{\mu} + (1-\hat{\mu})\psi
\end{equation}
これらを用いて運動方程式を整理すると：
\begin{equation}
  \rho(\psi)\!\left(\pd{\bu}{t} + \bu\cdot\bnabla\bu\right)
  = -\bnabla p
  + \frac{1}{Re}\bnabla\cdot\!\left[\mu(\psi)\!\left(\bnabla\bu + \bnabla\bu^T\right)\right]
  + \frac{\rho(\psi)}{Fr^2}\hat{\bg}
  + \frac{1}{We}\,\kappa\,\bnabla\psi
\end{equation}


%% ============================================================
\section{Conservative Level Set 法}
\label{sec:levelset}
%% ============================================================

\subsection{保存形移流方程式}

Conservative Level Set 法では、界面を $0$ から $1$ へ変化する滑らかな階段関数 $\psi$ で表現する。
$\psi=0.5$ が界面位置に対応する。この関数は流速 $\bu$ に沿って以下の保存形で移流される：
\begin{equation}
  \pd{\psi}{t} + \bnabla \cdot (\psi \bu) = 0
\end{equation}
非圧縮条件 $\bnabla \cdot \bu = 0$ を用いると非保存形 $\partial \psi / \partial t + \bu \cdot \bnabla \psi = 0$ と等価であるが、
数値計算上は保存形を維持することで全質量の保存性を高める。

\subsection{5次精度 WENO スキームによる離散化}

界面の鋭さを維持しつつ数値振動を抑えるため、移流項の離散化には \textbf{5次精度 WENO (Weighted Essentially Non-Oscillatory) スキーム} を使用する。
フラックス $f = \psi u$ に対して、セル界面での値を周囲のスタシルから適応的に補間する：
\begin{equation}
  \hat{f}_{i+1/2} = \sum_{k=1}^3 \omega_k \hat{f}^{(k)}_{i+1/2}
\end{equation}
ここで $\omega_k$ は解の滑らかさに応じて変化する重み係数である。

\subsection{再圧縮（Re-initialization）ステップ}

移流によって拡散・歪んだ $\psi$ プロファイルを、一定の厚み $\varepsilon$ を持つプロファイルへ維持するため、
各タイムステップ後に以下の修正方程式を定常状態まで解く：
\begin{equation}
  \pd{\psi}{\tau} + \bnabla \cdot [ \psi(1-\psi) \bm{n} ] = \varepsilon \bnabla \cdot [ (\bnabla \psi \cdot \bm{n}) \bm{n} ]
\end{equation}
ここで $\bm{n} = \bnabla \psi / |\bnabla \psi|$ は界面法線ベクトルである。
このステップにより、界面付近の $\psi$ は $\psi(\bm{x}) = \frac{1}{2}(1 + \tanh(d/2\varepsilon))$ 
（$d$ は界面からの距離）という安定した形状に保たれる。

\subsection{物性値の補間}

密度 $\rho$ および粘性 $\mu$ は，$\psi$ を用いて全領域で定義される：
\begin{align}
  \rho(\psi) &= \rho_g + (\rho_l - \rho_g)\psi \\
  \mu(\psi)  &= \mu_g + (\mu_l - \mu_g)\psi
\end{align}
これにより、気相 ($\psi=0$) から液相 ($\psi=1$) まで滑らかに物性が変化し、不連続性に起因する数値的不安定が解消される。


%% ============================================================
\section{界面適合格子の生成}
\label{sec:grid}
%% ============================================================

\subsection{グリッド密度関数による格子制御}

計算領域 $[0, L_x]\times[0, L_y]$ において，
界面（$\psi=0.5$）付近の解像度を局所的に高めるため，
一様計算空間 $(\xi, \eta) \in [0, 1]\times[0, 1]$ から
実空間 $(x, y)$ への写像 $x(\xi), y(\eta)$ を導入する。

グリッド密度関数 $\omega(\psi)$ を以下のように定義する：
\begin{equation}
  \omega(\psi) = 1 + \alpha \exp\left( - \beta (\psi - 0.5)^2 \right)
\end{equation}
ここで $\alpha, \beta$ は解像度を制御するパラメータである。
格子の間隔 $\Delta x$ は $1/\omega$ に比例するように配置される。


%% ============================================================
\section{結合コンパクト差分法 (CCD)}
\label{sec:CCD}
%% ============================================================

\subsection{CCD法の定式化}

格子点 $x_i$ における関数値を $f_i$, 一次微分値を $f'_i$, 二次微分値を $f''_i$ とし、
これらを同時に未知数として解くのが結合コンパクト差分法（CCD法）である。
Chu \& Fan \cite{ChuFan1998} に基づき、$\Ord{h^6}$ 精度の関係式を以下のように構成する。

\begin{tcolorbox}[algbox, title={CCD $\Ord{h^6}$ 差分式}]
一次微分 $f'$ と二次微分 $f''$ の間の連立関係式：
\begin{align}
  \alpha_1 (f'_{i+1} + f'_{i-1}) + f'_i + \beta_1 h (f''_{i+1} - f''_{i-1}) &= \frac{a_1}{2h} (f_{i+1} - f_{i-1}) \label{eq:ccd1} \\
  \alpha_2 \frac{1}{h} (f'_{i+1} - f'_{i-1}) + f''_i + \beta_2 (f''_{i+1} + f''_{i-1}) &= \frac{a_2}{h^2} (f_{i+1} - 2f_i + f_{i-1}) \label{eq:ccd2}
\end{align}
\end{tcolorbox}

テイラー展開による誤差解析から得られる $\Ord{h^6}$ 精度の係数は以下の通りである：
\begin{equation}
  \alpha_1 = \frac{7}{16},\, \beta_1 = -\frac{1}{16},\, a_1 = \frac{15}{8},\quad
  \alpha_2 = \frac{9}{8h},\, \beta_2 = \frac{1}{8},\, a_2 = \frac{3}{1} (\text{修正が必要な場合あり})
\end{equation}
(正確な係数導出の詳細は省略するが、これらは 3 重対角行列の反転によって効率的に解くことができる。)

\subsection{混合偏微分の計算}

曲率計算に必要な $f_{xy}$ は、まず $x$ 方向に一次微分を求め、その結果に対して $y$ 方向に
同様の CCD 一次微分を適用することで求める。


%% ============================================================
\section{コロケート格子と Rhie--Chow 補正}
\label{sec:collocate}
%% ============================================================

すべての変数（$\bu, p, \psi$）を格子点（セル中心）で定義するコロケート格子を採用する。
スタガード格子で発生する座標変換の複雑さを避けるためであるが、
そのままでは圧力のチェッカーボード不安定が発生する。

これを防ぐため，セル界面での流速 $U_{i+1/2}$ を計算する際に，
Rhie--Chow 型の補正項（圧力の三階微分に相当する散逸）を付加する：
\begin{equation}
  U_{i+1/2} = \overline{u}_{i+1/2} - C \Delta x^2 \left( \frac{\partial^3 p}{\partial x^3} \right)_{i+1/2}
\end{equation}
この界面流速を用いて連続の式（圧力 Poisson 方程式のソース項）を構成する。


%% ============================================================
\section{圧力 Poisson 方程式の求解}
\label{sec:pressure}
%% ============================================================

\subsection{変密度 Poisson 方程式の導出}

Navier--Stokes 方程式を時間離散化し、連続の式 $\bnabla \cdot \bu^{n+1} = 0$ を要請すると、
圧力に関する以下の楕円型方程式が得られる：
\begin{equation}
  \bnabla \cdot \left( \frac{1}{\rho(\psi)} \bnabla p^{n+1} \right) = \frac{\bnabla \cdot \bu^*}{\Delta t}
\end{equation}
ここで $\bu^*$ は圧力項を除いた運動量予測ステップで得られる中間速度である。

\subsection{数値解法}

この方程式を CCD 法により空間離散化すると、非対称な大型連立一次方程式となる。
これを BiCGSTAB 法（前処理付き共役勾配法の一種）を用いて反復求解する。
収束判定基準は残差ノルムが $10^{-8}$ 以下とする。


%% ============================================================
\section{時間積分スキーム}
\label{sec:time}
%% ============================================================

移流項（速度および Level Set）には、安定性の高い \textbf{3次精度 TVD Runge--Kutta 法} を用いる：
\begin{align}
  \phi^{(1)} &= \phi^n + \Delta t \mathcal{L}(\phi^n) \notag \\
  \phi^{(2)} &= \frac{3}{4}\phi^n + \frac{1}{4}\phi^{(1)} + \frac{1}{4}\Delta t \mathcal{L}(\phi^{(1)}) \notag \\
  \phi^{n+1} &= \frac{1}{3}\phi^n + \frac{2}{3}\phi^{(2)} + \frac{2}{3}\Delta t \mathcal{L}(\phi^{(2)})
\end{align}
ここで $\mathcal{L}$ は WENO5 による移流演算子である。
一方、粘性項については、拡散数制限を緩和するため Crank--Nicolson 法による半陰的処理を行う。


%% ============================================================
\section{完全アルゴリズム}
\label{sec:algorithm}
%% ============================================================

各タイムステップ $t^n \to t^n + \Delta t$ における計算手順を以下にまとめる。

\begin{tcolorbox}[algbox, title={計算アルゴリズムフロー}]
\begin{enumerate}
  \item \textbf{Conservative Level Set の移流}：WENO5 と TVD-RK3 を用い、$\psi^n$ を移流して $\psi^*$ を得る。
  \item \textbf{再圧縮ステップ}：$\psi^*$ に対して修正方程式を数ステップ解き、プロファイルを維持した $\psi^{n+1}$ を決定する。
  \item \textbf{物性値と曲率の更新}：$\psi^{n+1}$ から $\rho, \mu$ を求め、CCD法で曲率 $\kappa^{n+1}$ を計算する。
  \item \textbf{運動量予測（Predictor）}：圧力項を除いた NS 方程式を解き、中間速度 $\bu^*$ を求める。
  \item \textbf{圧力 Poisson 方程式の求解}：変密度 PPE を BiCGSTAB で解き、圧力 $p^{n+1}$ を決定する。
  \item \textbf{速度補正（Corrector）}：$\bu^*$, $p^{n+1}$, $\bnabla p^{n+1}$（CCD）および Rhie--Chow 補正を用いて、
        非圧縮条件を満たす速度 $\bu^{n+1}$ を計算する。
  \item \textbf{時間刻みの更新}：CFL 条件に基づき、次ステップの $\Delta t$ を決定する。
\end{enumerate}
\end{tcolorbox}


%% ============================================================
\begin{thebibliography}{99}
\bibitem[Brackbill et al.(1992)]{Brackbill1992}
Brackbill, J.~U., Kothe, D.~B., \& Zemach, C. (1992).
A Continuum Method for Modeling Surface Tension.
\textit{Journal of Computational Physics}, \textbf{100}(2), 335--354.

\bibitem[Chorin(1968)]{Chorin1968}
Chorin, A.~J. (1968).
Numerical Solution of the Navier--Stokes Equations.
\textit{Mathematics of Computation}, \textbf{22}(104), 745--762.

\bibitem[Chu \& Fan(1998)]{ChuFan1998}
Chu, K.~W., \& Fan, L.~S. (1998).
A Combined Compact Difference Scheme for Numerical Simulation of Particle-Fluid Flows.
\textit{Journal of Computational Physics}, \textbf{141}(1), 1--24.

\bibitem[Olsson \& Kreiss(2005)]{Olsson2005}
Olsson, E., \& Kreiss, G. (2005).
A Conservative Level Set Method for Two Phase Flow.
\textit{Journal of Computational Physics}, \textbf{210}(1), 225--246.

\bibitem[Shu \& Osher(1988)]{Shu1988}
Shu, C.-W., \& Osher, S. (1988).
Efficient Implementation of Essentially Non-Oscillatory Shock-Capturing Schemes.
\textit{Journal of Computational Physics}, \textbf{77}(2), 439--471.

\bibitem[Sussman et al.(1994)]{Sussman1994}
Sussman, M., Smereka, P., \& Osher, S. (1994).
A Level Set Approach for Computing Solutions to Incompressible Two-Phase Flow.
\textit{Journal of Computational Physics}, \textbf{114}(1), 146--159.

\bibitem[Unverdi \& Tryggvason(1992)]{Unverdi1992}
Unverdi, S.~O., \& Tryggvason, G. (1992).
A Front-Tracking Method for Viscous, Incompressible, Multi-fluid Flows.
\textit{Journal of Computational Physics}, \textbf{100}(1), 25--37.
\end{thebibliography}

\end{document}